
\mxunit{stmg-c01}{Suites numériques}
\mxskills{chercher, modéliser, représenter, calculer, communiquer}

\section{Cours de référence}
Une suite numérique associe à chaque entier naturel $n$ un nombre $u_n$. L'indice $n$ désigne un rang ; il ne doit pas être confondu avec la valeur $u_n$.

\subsection{Suite arithmétique}
\begin{keypoint}
Une suite $(u_n)$ est \textbf{arithmétique} lorsqu'il existe un réel $r$ tel que, pour tout $n$, $u_{n+1}=u_n+r$. Le nombre $r$ est la raison. Si le premier terme est $u_0$, alors $u_n=u_0+nr$. Si le premier terme est $u_1$, alors $u_n=u_1+(n-1)r$.
\end{keypoint}

\begin{center}
\begin{tikzpicture}[>=Stealth,node distance=18mm]
\node[draw,rounded corners,fill=mxpaper] (u0) {$u_0$};
\node[draw,rounded corners,fill=mxpaper,right=of u0] (u1) {$u_1$};
\node[draw,rounded corners,fill=mxpaper,right=of u1] (u2) {$u_2$};
\node[draw,rounded corners,fill=mxpaper,right=of u2] (un) {$u_n$};
\draw[->,thick] (u0)--node[above] {$+r$}(u1);
\draw[->,thick] (u1)--node[above] {$+r$}(u2);
\draw[->,thick,dashed] (u2)--node[above] {$+(n-2)r$}(un);
\end{tikzpicture}
\end{center}

Trois nombres $x,y,z$, écrits dans cet ordre, sont trois termes consécutifs d'une suite arithmétique si $y-x=z-y$, c'est-à-dire si $y=\dfrac{x+z}{2}$. Cette dernière quantité est la moyenne arithmétique de $x$ et $z$.

\begin{method}
Pour modéliser une situation additive : 1) définir $u_n$ et préciser l'unité ; 2) identifier le premier terme ; 3) traduire l'augmentation ou la diminution constante par la raison $r$ ; 4) écrire le terme général ; 5) répondre par une phrase contextualisée.
\end{method}

\begin{example}
Un commerce vend $480$ unités le premier mois puis $25$ unités de plus chaque mois. Avec $u_1=480$ et $r=25$, on obtient $u_n=480+25(n-1)$. Ainsi $u_{12}=755$ : le commerce prévoit $755$ unités vendues au douzième mois.
\end{example}

\subsection{Somme de termes arithmétiques}
Pour $n+1$ termes de $u_0$ à $u_n$,
\[
u_0+u_1+\cdots+u_n=(n+1)\frac{u_0+u_n}{2}.
\]
La formule se mémorise sous la forme : \emph{nombre de termes multiplié par la moyenne du premier et du dernier terme}.

\begin{pitfall}
Dans $u_3+u_4+\cdots+u_{12}$, il y a $12-3+1=10$ termes, et non $9$ ni $12$.
\end{pitfall}

\subsection{Suite géométrique}
\begin{keypoint}
Une suite $(u_n)$ est \textbf{géométrique} lorsqu'il existe un réel $q$ tel que $u_{n+1}=q u_n$. Si le premier terme est $u_0$, alors $u_n=u_0q^n$. Si le premier terme est $u_1$, alors $u_n=u_1q^{n-1}$.
\end{keypoint}

\begin{center}
\begin{tikzpicture}[>=Stealth,node distance=18mm]
\node[draw,rounded corners,fill=mxmintpale] (u0) {$u_0$};
\node[draw,rounded corners,fill=mxmintpale,right=of u0] (u1) {$u_1$};
\node[draw,rounded corners,fill=mxmintpale,right=of u1] (u2) {$u_2$};
\node[draw,rounded corners,fill=mxmintpale,right=of u2] (un) {$u_n$};
\draw[->,thick] (u0)--node[above] {$\times q$}(u1);
\draw[->,thick] (u1)--node[above] {$\times q$}(u2);
\draw[->,thick,dashed] (u2)--node[above] {$\times q^{n-2}$}(un);
\end{tikzpicture}
\end{center}

Une hausse de $t\,\%$ correspond à $q=1+\dfrac{t}{100}$ ; une baisse de $t\,\%$ à $q=1-\dfrac{t}{100}$. Trois termes strictement positifs $x,y,z$ sont consécutifs dans une suite géométrique si $y^2=xz$ ; $y=\sqrt{xz}$ est leur moyenne géométrique.

Pour $q\neq1$,
\[
u_0+u_1+\cdots+u_n=u_0\frac{1-q^{n+1}}{1-q}.
\]

\subsection{Algorithmique : calculer une somme}
\begin{lstlisting}[style=python]
def somme_geometrique(n):
    s = 0
    for k in range(n + 1):
        s = s + 5 * 2**k
    return s
\end{lstlisting}
L'accumulateur \texttt{s} contient la somme partielle ; le compteur \texttt{k} parcourt les entiers de $0$ à $n$. Ici \texttt{somme\_geometrique(4)} calcule $5+10+20+40+80=155$, valeur que l'on retrouve avec $u_0\frac{1-q^{n+1}}{1-q}=5\times\frac{1-2^5}{1-2}$ pour la suite géométrique de premier terme $5$ et de raison $2$.

\mxsession{1}{\mxdate*{2026-09-15}}{B : 2 h -- C : 1 h}{Suite arithmétique : reconnaître et modéliser}{Suites numériques}
    \begin{objectives}\begin{itemize}\item Reconnaître une évolution additive.\item Déterminer la raison.\item Définir une suite dans un contexte.\end{itemize}\end{objectives}
    \begin{fluency}[6 à 8 minutes, sans calculatrice]\begin{enumerate}\item Écrire le coefficient multiplicateur d'une hausse de $\pc{10}$. \mxblank{$1+\dfrac{10}{100}=\num{1,10}$}\item Calculer $u_4$ si $u_0=14$ et $u_{n+1}=u_n+3$. \mxblank{$u_4=26$}\item Combien de termes de $u_3$ à $u_{15}$ ? \mxblank{$15-3+1=13$}\item Donner le coefficient d'une baisse de $\pc{8}$. \mxblank{$\num{0,92}$}\item Donner la raison de $12\ ;\ 12+4\ ;\ 12+2\times4\ ;\ \ldots$. \mxblank{$r=4$}\end{enumerate}\end{fluency}
    \begin{prerequisites}Relire la dernière trace écrite et savoir effectuer les calculs élémentaires mobilisés aujourd'hui. Avant l'activité, faire reformuler les mots-clés : reconnaître une évolution additive et déterminer la raison.\end{prerequisites}

    \begin{activity}
    Une association compte $320$ adhérents. Elle prévoit d'en gagner $18$ chaque année. On note $u_n$ l'effectif après $n$ années. Calculer $u_0,u_1,u_2$. Quelle opération fait passer d'un terme au suivant ? Écrire une relation entre $u_{n+1}$ et $u_n$, puis prévoir $u_5$.
    \begin{center}\begin{tikzpicture}[>=Stealth,node distance=17mm]
\foreach \i in {0,...,4}{\node[draw,rounded corners,fill=mxpaper] (u\i) at (1.7*\i,0) {$u_\i$};}
\foreach \i [evaluate=\i as \j using int(\i+1)] in {0,...,3}{\draw[->,thick] (u\i)--node[above] {$ +r $}(u\j);}
\end{tikzpicture}\end{center}
    \end{activity}
    \begin{solution}$u_0=320$, $u_1=338$, $u_2=356$ ; on ajoute $18$. Donc $u_{n+1}=u_n+18$ et $u_5=410$. Faire verbaliser la donnée qui déclenche le choix de méthode.\end{solution}

    \begin{summary}Une suite est arithmétique lorsque $u_{n+1}-u_n=r$ est constant. Dans une situation, une augmentation absolue constante conduit à une suite arithmétique.\end{summary}

    \begin{example}
    Reprendre les données de l'activité et rédiger la solution dans l'ordre : définition des objets, propriété utilisée, calcul exact ou approché, puis phrase de réponse. La réponse modèle figure dans l'édition enseignant.
    \end{example}
    \begin{solution}$u_0=320$, $u_1=338$, $u_2=356$ ; on ajoute $18$. Donc $u_{n+1}=u_n+18$ et $u_5=410$.\end{solution}

    \begin{exercise}[essentiel]
    \begin{enumerate}\item La suite $(u_n)$ vérifie $u_0=110$ et $u_{n+1}=u_n+3$. Écrire $u_n$ puis calculer $u_{12}$.\item Une quantité vaut $420$ et baisse de $\pc{3}$ par période. Écrire $v_n$ et calculer $v_5$.\item Pour chacune des suites $(u_n)$ et $(v_n)$ ci-dessus, préciser si le modèle est arithmétique ou géométrique et justifier par une différence ou un quotient.\end{enumerate}
    \end{exercise}
    \begin{solution}$u_n=110+3n$ et $u_{12}=146$. Pour la seconde, $q=\num{0,97}$, $v_n=420\times\num{0,97}^n$ et $v_5\approx\num{360,67}$. La suite $(u_n)$ est arithmétique (différence constante $3$) et $(v_n)$ est géométrique (quotient constant $\num{0,97}$).\end{solution}

\mxexuse{1}

    \begin{homework}Reprendre la question 2 de l'exercice essentiel de la séance avec les données suivantes : remplacer la donnée initiale par $125$ et le taux ou la raison par $3$. Rédiger en quatre lignes : données, formule, calcul, conclusion.\end{homework}
    \begin{solution}La correction dépend de l'exercice repris. Vérifier surtout la cohérence du modèle, le nombre d'itérations, l'arrondi et la phrase de conclusion.\end{solution}

    \begin{exitticket}[3 minutes]Sur une demi-feuille : écrire la règle principale de la séance, traiter un exemple très court et indiquer son degré de confiance de 1 à 4. Question(s) essentielle(s) du jour : 1, 2 et 3.\end{exitticket}

    \ifmxteacher
    \begin{method}
    \textbf{Conduite en 1 h :} 7 min d'automatismes ; 10 min de mise en commun et trace ; 35 min d'exercices essentiels ; 5 min de correction ; 3 min de billet de sortie.\\
    \textbf{Conduite en 2 h :} ajouter 20 min de recherche de l'activité, 20 min d'institutionnalisation, un second exercice autonome et 10 min de correction comparée.\\
    \textbf{Différenciation :} fournir aux élèves fragiles une fiche à quatre cases (données, modèle, calcul, conclusion). Demander aux élèves rapides de produire un contre-exemple, une justification ou une seconde méthode.
    \end{method}
    \fi
\mxsession{2}{\mxdate*{2026-09-17}}{B : 1 h -- C : 2 h}{Suite arithmétique : terme général}{Suites numériques}
\begin{objectives}\begin{itemize}\item Passer de la récurrence au terme général.\item Calculer un terme éloigné.\item Interpréter le rang.\end{itemize}\end{objectives}
\begin{fluency}[6 à 8 minutes, sans calculatrice]\begin{enumerate}\item Réduire $3x+5-x+2$. \mxblank{$2x+7$}\item Combien de termes de $u_{2}$ à $u_{10}$ ? \mxblank{$10-2+1=9$}\item Donner le coefficient d'une baisse de $\pc{12}$. \mxblank{$\num{0,88}$}\item Donner la raison de $12\ ;\ 12+5\ ;\ 12+2\times5\ ;\ \ldots$. \mxblank{$r=5$}\item Calculer $u_4$ si $u_0=15$ et $u_{n+1}=u_n+3$. \mxblank{$u_4=27$}\end{enumerate}\end{fluency}
\begin{prerequisites}Relire la dernière trace écrite et savoir effectuer les calculs élémentaires mobilisés aujourd'hui. Avant l'activité, faire reformuler les mots-clés : passer de la récurrence au terme général et calculer un terme éloigné.\end{prerequisites}

\begin{activity}
Un magasin possède $\num{1250}$ cartes de fidélité actives au rang $0$ et en gagne $75$ par mois. Sans calculer tous les termes, déterminer l'effectif au rang $n$, puis au rang $18$. Expliquer le rôle de $18$ dans le calcul.

\end{activity}
\begin{solution}$u_n=\num{1250}+75n$, donc $u_{18}=\num{2600}$. Le nombre $18$ compte les augmentations depuis le rang $0$. Faire verbaliser la donnée qui déclenche le choix de méthode.\end{solution}

\begin{summary}Si $(u_n)$ est arithmétique de raison $r$, alors $u_n=u_0+nr$ ou $u_n=u_1+(n-1)r$, selon le rang du premier terme.\end{summary}

\begin{example}
Reprendre les données de l'activité et rédiger la solution dans l'ordre : définition des objets, propriété utilisée, calcul exact ou approché, puis phrase de réponse. La réponse modèle figure dans l'édition enseignant.
\end{example}
\begin{solution}$u_n=\num{1250}+75n$, donc $u_{18}=\num{2600}$. Le nombre $18$ compte les augmentations depuis le rang $0$.\end{solution}

\begin{exercise}[essentiel]
\begin{enumerate}\item La suite $(u_n)$ vérifie $u_0=120$ et $u_{n+1}=u_n+4$. Écrire $u_n$ puis calculer $u_{12}$.\item Une quantité vaut $440$ et baisse de $\pc{4}$ par période. Écrire $v_n$ et calculer $v_5$.\item Pour chacune des suites $(u_n)$ et $(v_n)$ ci-dessus, préciser si le modèle est arithmétique ou géométrique et justifier par une différence ou un quotient.\end{enumerate}
\end{exercise}
\begin{solution}$u_n=120+4n$ et $u_{12}=168$. Pour la seconde, $q=\num{0,96}$, $v_n=440\times\num{0,96}^n$ et $v_5\approx\num{358,76}$. La suite $(u_n)$ est arithmétique (différence constante $4$) et $(v_n)$ est géométrique (quotient constant $\num{0,96}$).\end{solution}

\mxexuse{2}

\begin{homework}Reprendre la question 2 de l'exercice essentiel de la séance avec les données suivantes : remplacer la donnée initiale par $130$ et le taux ou la raison par $4$. Rédiger en quatre lignes : données, formule, calcul, conclusion.\end{homework}
\begin{solution}La correction dépend de l'exercice repris. Vérifier surtout la cohérence du modèle, le nombre d'itérations, l'arrondi et la phrase de conclusion.\end{solution}

\begin{exitticket}[3 minutes]Sur une demi-feuille : écrire la règle principale de la séance, traiter un exemple très court et indiquer son degré de confiance de 1 à 4. Question(s) essentielle(s) du jour : 1 et 2.\end{exitticket}

\ifmxteacher
\begin{method}
\textbf{Conduite en 1 h :} 7 min d'automatismes ; 10 min de mise en commun et trace ; 35 min d'exercices essentiels ; 5 min de correction ; 3 min de billet de sortie.\\
\textbf{Conduite en 2 h :} ajouter 20 min de recherche de l'activité, 20 min d'institutionnalisation, un second exercice autonome et 10 min de correction comparée.\\
\textbf{Différenciation :} fournir aux élèves fragiles une fiche à quatre cases (données, modèle, calcul, conclusion). Demander aux élèves rapides de produire un contre-exemple, une justification ou une seconde méthode.
\end{method}
\fi
\mxsession{3}{\mxdate*{2026-09-22}}{B : 2 h -- C : 1 h}{Somme des termes d'une suite arithmétique}{Suites numériques}
\begin{objectives}\begin{itemize}\item Compter les termes d'une somme.\item Appliquer la formule.\item Contrôler l'ordre de grandeur.\end{itemize}\end{objectives}
\begin{fluency}[6 à 8 minutes, sans calculatrice]\begin{enumerate}\item Calculer $\pc{20}$ de $120$. \mxblank{$24$}\item Donner le coefficient d'une baisse de $\pc{15}$. \mxblank{$\num{0,85}$}\item Donner la raison de $12\ ;\ 12+6\ ;\ 12+2\times6\ ;\ \ldots$. \mxblank{$r=6$}\item Calculer $u_4$ si $u_0=16$ et $u_{n+1}=u_n+3$. \mxblank{$u_4=28$}\item Combien de termes de $u_{5}$ à $u_{20}$ ? \mxblank{$20-5+1=16$}\end{enumerate}\end{fluency}
\begin{prerequisites}Relire la dernière trace écrite et savoir effectuer les calculs élémentaires mobilisés aujourd'hui. Avant l'activité, faire reformuler les mots-clés : compter les termes d'une somme et appliquer la formule.\end{prerequisites}

\begin{activity}
Une entreprise verse une prime mensuelle : $45$ euros le premier mois, puis $5$ euros de plus chaque mois. Calculer le total versé pendant les douze premiers mois de trois façons : somme explicite, regroupement premier-dernier, formule.

\end{activity}
\begin{solution}$u_{12}=45+11\times5=100$. Le total vaut $12\times\dfrac{45+100}{2}=870$ euros. Faire verbaliser la donnée qui déclenche le choix de méthode.\end{solution}

\begin{summary}Une somme de termes arithmétiques vaut $\text{nombre de termes}\times\dfrac{\text{premier}+\text{dernier}}2$. De $u_p$ à $u_q$, il y a $q-p+1$ termes.\end{summary}

\begin{example}
Reprendre les données de l'activité et rédiger la solution dans l'ordre : définition des objets, propriété utilisée, calcul exact ou approché, puis phrase de réponse. La réponse modèle figure dans l'édition enseignant.
\end{example}
\begin{solution}$u_{12}=45+11\times5=100$. Le total vaut $12\times\dfrac{45+100}{2}=870$ euros.\end{solution}

\begin{exercise}[essentiel]
\begin{enumerate}\item La suite $(u_n)$ vérifie $u_0=130$ et $u_{n+1}=u_n+5$. Écrire $u_n$ puis calculer $u_{12}$.\item Une quantité vaut $460$ et baisse de $\pc{5}$ par période. Écrire $v_n$ et calculer $v_5$.\item Pour chacune des suites $(u_n)$ et $(v_n)$ ci-dessus, préciser si le modèle est arithmétique ou géométrique et justifier par une différence ou un quotient.\end{enumerate}
\end{exercise}
\begin{solution}$u_n=130+5n$ et $u_{12}=190$. Pour la seconde, $q=\num{0,95}$, $v_n=460\times\num{0,95}^n$ et $v_5\approx\num{355,94}$. La suite $(u_n)$ est arithmétique (différence constante $5$) et $(v_n)$ est géométrique (quotient constant $\num{0,95}$).\end{solution}

\mxexuse{3}

\begin{homework}Reprendre la question 2 de l'exercice essentiel de la séance avec les données suivantes : remplacer la donnée initiale par $135$ et le taux ou la raison par $5$. Rédiger en quatre lignes : données, formule, calcul, conclusion.\end{homework}
\begin{solution}La correction dépend de l'exercice repris. Vérifier surtout la cohérence du modèle, le nombre d'itérations, l'arrondi et la phrase de conclusion.\end{solution}

\begin{exitticket}[3 minutes]Sur une demi-feuille : écrire la règle principale de la séance, traiter un exemple très court et indiquer son degré de confiance de 1 à 4. Question(s) essentielle(s) du jour : 1, 2 et 3.\end{exitticket}

\ifmxteacher
\begin{method}
\textbf{Conduite en 1 h :} 7 min d'automatismes ; 10 min de mise en commun et trace ; 35 min d'exercices essentiels ; 5 min de correction ; 3 min de billet de sortie.\\
\textbf{Conduite en 2 h :} ajouter 20 min de recherche de l'activité, 20 min d'institutionnalisation, un second exercice autonome et 10 min de correction comparée.\\
\textbf{Différenciation :} fournir aux élèves fragiles une fiche à quatre cases (données, modèle, calcul, conclusion). Demander aux élèves rapides de produire un contre-exemple, une justification ou une seconde méthode.
\end{method}
\fi
\mxsession{4}{\mxdate*{2026-09-24}}{B : 1 h -- C : 2 h}{Suite géométrique : reconnaître et modéliser}{Suites numériques}
    \begin{objectives}\begin{itemize}\item Reconnaître une évolution multiplicative.\item Passer d'un taux à un coefficient.\item Comparer arithmétique et géométrique.\end{itemize}\end{objectives}
    \begin{fluency}[6 à 8 minutes, sans calculatrice]\begin{enumerate}\item Écrire le coefficient multiplicateur d'une hausse de $\pc{25}$. \mxblank{$1+\dfrac{25}{100}=\num{1,25}$}\item Donner la raison de $12\ ;\ 12+7\ ;\ 12+2\times7\ ;\ \ldots$. \mxblank{$r=7$}\item Calculer $u_4$ si $u_0=17$ et $u_{n+1}=u_n+3$. \mxblank{$u_4=29$}\item Combien de termes de $u_{0}$ à $u_{9}$ ? \mxblank{$9-0+1=10$}\item Donner le coefficient d'une baisse de $\pc{5}$. \mxblank{$\num{0,95}$}\end{enumerate}\end{fluency}
    \begin{prerequisites}Relire la dernière trace écrite et savoir effectuer les calculs élémentaires mobilisés aujourd'hui. Avant l'activité, faire reformuler les mots-clés : reconnaître une évolution multiplicative et passer d'un taux à un coefficient.\end{prerequisites}

    \begin{activity}
    Une boutique en ligne augmente son nombre de visiteurs de $\pc{6}$ par mois à partir de $\num{8000}$ visiteurs. Calculer les trois premières valeurs. Comparer les différences puis les quotients. Identifier la nature de la suite.
    \begin{center}\begin{tikzpicture}[>=Stealth,node distance=17mm]
\foreach \i in {0,...,4}{\node[draw,rounded corners,fill=mxpaper] (u\i) at (1.7*\i,0) {$u_\i$};}
\foreach \i [evaluate=\i as \j using int(\i+1)] in {0,...,3}{\draw[->,thick] (u\i)--node[above] {$ \times q $}(u\j);}
\end{tikzpicture}\end{center}
    \end{activity}
    \begin{solution}Le coefficient est $\num{1,06}$ : $\num{8000}$, $\num{8480}$, $\num{8988,8}$. Les quotients valent $\num{1,06}$ ; la suite est géométrique. Faire verbaliser la donnée qui déclenche le choix de méthode.\end{solution}

    \begin{summary}Une suite est géométrique lorsque $\dfrac{u_{n+1}}{u_n}=q$ est constant (termes non nuls). Une évolution relative constante conduit à une suite géométrique.\end{summary}

    \begin{example}
    Reprendre les données de l'activité et rédiger la solution dans l'ordre : définition des objets, propriété utilisée, calcul exact ou approché, puis phrase de réponse. La réponse modèle figure dans l'édition enseignant.
    \end{example}
    \begin{solution}Le coefficient est $\num{1,06}$ : $\num{8000}$, $\num{8480}$, $\num{8988,8}$. Les quotients valent $\num{1,06}$ ; la suite est géométrique.\end{solution}

    \begin{exercise}[essentiel]
    \begin{enumerate}\item La suite $(u_n)$ vérifie $u_0=140$ et $u_{n+1}=u_n+6$. Écrire $u_n$ puis calculer $u_{12}$.\item Une quantité vaut $480$ et baisse de $\pc{6}$ par période. Écrire $v_n$ et calculer $v_5$.\item Pour chacune des suites $(u_n)$ et $(v_n)$ ci-dessus, préciser si le modèle est arithmétique ou géométrique et justifier par une différence ou un quotient.\end{enumerate}
    \end{exercise}
    \begin{solution}$u_n=140+6n$ et $u_{12}=212$. Pour la seconde, $q=\num{0,94}$, $v_n=480\times\num{0,94}^n$ et $v_5\approx\num{352,27}$. La suite $(u_n)$ est arithmétique (différence constante $6$) et $(v_n)$ est géométrique (quotient constant $\num{0,94}$).\end{solution}

\mxexuse{4}

    \begin{homework}Reprendre la question 2 de l'exercice essentiel de la séance avec les données suivantes : remplacer la donnée initiale par $140$ et le taux ou la raison par $6$. Rédiger en quatre lignes : données, formule, calcul, conclusion.\end{homework}
    \begin{solution}La correction dépend de l'exercice repris. Vérifier surtout la cohérence du modèle, le nombre d'itérations, l'arrondi et la phrase de conclusion.\end{solution}

    \begin{exitticket}[3 minutes]Sur une demi-feuille : écrire la règle principale de la séance, traiter un exemple très court et indiquer son degré de confiance de 1 à 4. Question(s) essentielle(s) du jour : 1 et 2.\end{exitticket}

    \ifmxteacher
    \begin{method}
    \textbf{Conduite en 1 h :} 7 min d'automatismes ; 10 min de mise en commun et trace ; 35 min d'exercices essentiels ; 5 min de correction ; 3 min de billet de sortie.\\
    \textbf{Conduite en 2 h :} ajouter 20 min de recherche de l'activité, 20 min d'institutionnalisation, un second exercice autonome et 10 min de correction comparée.\\
    \textbf{Différenciation :} fournir aux élèves fragiles une fiche à quatre cases (données, modèle, calcul, conclusion). Demander aux élèves rapides de produire un contre-exemple, une justification ou une seconde méthode.
    \end{method}
    \fi
\mxsession{5}{\mxdate*{2026-09-29}}{B : 2 h -- C : 1 h}{Suite géométrique : terme général et somme}{Suites numériques}
\begin{objectives}\begin{itemize}\item Écrire le terme général.\item Calculer une somme géométrique.\item Interpréter une évolution cumulée.\end{itemize}\end{objectives}
\begin{fluency}[6 à 8 minutes, sans calculatrice]\begin{enumerate}\item Réduire $3x+8-x+2$. \mxblank{$2x+10$}\item Calculer $u_4$ si $u_0=18$ et $u_{n+1}=u_n+3$. \mxblank{$u_4=30$}\item Combien de termes de $u_{4}$ à $u_{11}$ ? \mxblank{$11-4+1=8$}\item Donner le coefficient d'une baisse de $\pc{30}$. \mxblank{$\num{0,70}$}\item Donner la raison de $12\ ;\ 12+8\ ;\ 12+2\times8\ ;\ \ldots$. \mxblank{$r=8$}\end{enumerate}\end{fluency}
\begin{prerequisites}Relire la dernière trace écrite et savoir effectuer les calculs élémentaires mobilisés aujourd'hui. Avant l'activité, faire reformuler les mots-clés : écrire le terme général et calculer une somme géométrique.\end{prerequisites}

\begin{activity}
Un capital de $\num{250000}$ Ar est placé à $\pc{4}$ par an, intérêts composés. Écrire le capital $C_n$ après $n$ années. Calculer $C_6$, puis la somme $C_0+\cdots+C_6$ et expliquer ce que cette dernière somme représente mathématiquement et pourquoi elle n'est pas le capital final.

\end{activity}
\begin{solution}$C_n=\num{250000}\times\num{1,04}^n$ et $C_6\approx\num{316330}$ Ar. La somme $C_0+\cdots+C_6=\num{250000}\times\dfrac{1-\num{1,04}^7}{1-\num{1,04}}\approx\num{1974574}$ Ar est la somme de sept états successifs du capital, pas une accumulation réelle du placement. Faire verbaliser la donnée qui déclenche le choix de méthode.\end{solution}

\begin{summary}Si $(u_n)$ est géométrique de raison $q$, $u_n=u_0q^n$ ou $u_n=u_1q^{n-1}$. Pour $q\neq1$, $u_0+\cdots+u_n=u_0\dfrac{1-q^{n+1}}{1-q}$.\end{summary}

\begin{example}
Reprendre les données de l'activité et rédiger la solution dans l'ordre : définition des objets, propriété utilisée, calcul exact ou approché, puis phrase de réponse. La réponse modèle figure dans l'édition enseignant.
\end{example}
\begin{solution}$C_n=\num{250000}\times\num{1,04}^n$ et $C_6\approx\num{316330}$ Ar. La somme $C_0+\cdots+C_6=\num{250000}\times\dfrac{1-\num{1,04}^7}{1-\num{1,04}}\approx\num{1974574}$ Ar est la somme de sept états successifs du capital, pas une accumulation réelle du placement.\end{solution}

\begin{exercise}[essentiel]
\begin{enumerate}\item La suite $(u_n)$ vérifie $u_0=150$ et $u_{n+1}=u_n+2$. Écrire $u_n$ puis calculer $u_{12}$.\item Une quantité vaut $500$ et baisse de $\pc{2}$ par période. Écrire $v_n$ et calculer $v_5$.\item Pour chacune des suites $(u_n)$ et $(v_n)$ ci-dessus, préciser si le modèle est arithmétique ou géométrique et justifier par une différence ou un quotient.\end{enumerate}
\end{exercise}
\begin{solution}$u_n=150+2n$ et $u_{12}=174$. Pour la seconde, $q=\num{0,98}$, $v_n=500\times\num{0,98}^n$ et $v_5\approx\num{451,96}$. La suite $(u_n)$ est arithmétique (différence constante $2$) et $(v_n)$ est géométrique (quotient constant $\num{0,98}$).\end{solution}

\mxexuse{5}

\begin{homework}Reprendre la question 2 de l'exercice essentiel de la séance avec les données suivantes : remplacer la donnée initiale par $145$ et le taux ou la raison par $7$. Rédiger en quatre lignes : données, formule, calcul, conclusion.\end{homework}
\begin{solution}La correction dépend de l'exercice repris. Vérifier surtout la cohérence du modèle, le nombre d'itérations, l'arrondi et la phrase de conclusion.\end{solution}

\begin{exitticket}[3 minutes]Sur une demi-feuille : écrire la règle principale de la séance, traiter un exemple très court et indiquer son degré de confiance de 1 à 4. Question(s) essentielle(s) du jour : 1, 2 et 3.\end{exitticket}

\ifmxteacher
\begin{method}
\textbf{Conduite en 1 h :} 7 min d'automatismes ; 10 min de mise en commun et trace ; 35 min d'exercices essentiels ; 5 min de correction ; 3 min de billet de sortie.\\
\textbf{Conduite en 2 h :} ajouter 20 min de recherche de l'activité, 20 min d'institutionnalisation, un second exercice autonome et 10 min de correction comparée.\\
\textbf{Différenciation :} fournir aux élèves fragiles une fiche à quatre cases (données, modèle, calcul, conclusion). Demander aux élèves rapides de produire un contre-exemple, une justification ou une seconde méthode.
\end{method}
\fi
\mxsession{6}{\mxdate*{2026-10-01}}{B : 1 h -- C : 2 h}{Bilan et évaluation formative}{Suites numériques}
\begin{objectives}\begin{itemize}\item Choisir le bon modèle.\item Mobiliser les formules.\item Analyser ses erreurs.\end{itemize}\end{objectives}
\begin{fluency}[6 à 8 minutes, sans calculatrice]\begin{enumerate}\item Calculer $\pc{5}$ de $180$. \mxblank{$9$}\item Combien de termes de $u_{7}$ à $u_{18}$ ? \mxblank{$18-7+1=12$}\item Donner le coefficient d'une baisse de $\pc{20}$. \mxblank{$\num{0,80}$}\item Donner la raison de $12\ ;\ 12+9\ ;\ 12+2\times9\ ;\ \ldots$. \mxblank{$r=9$}\item Calculer $u_4$ si $u_0=19$ et $u_{n+1}=u_n+3$. \mxblank{$u_4=31$}\end{enumerate}\end{fluency}
\begin{prerequisites}Relire la dernière trace écrite et savoir effectuer les calculs élémentaires mobilisés aujourd'hui. Avant l'activité, faire reformuler les mots-clés : choisir le bon modèle et mobiliser les formules.\end{prerequisites}

\begin{activity}
Classer huit situations selon trois catégories : arithmétique, géométrique, ni l'une ni l'autre. Pour chaque situation, souligner le mot ou la donnée qui justifie le choix, puis choisir la formule adaptée.
\begin{enumerate}[label=\alph*)]
\item Un loyer de $\num{600000}$ Ar augmente de $\num{20000}$ Ar chaque année.
\item Une ville de $\num{5000}$ habitants voit sa population augmenter de $\pc{3}$ par an.
\item Un abonnement coûte $15$ euros à l'inscription, puis $2$ euros par séance : on note $u_n$ le coût total après $n$ séances.
\item Un véhicule perd $\pc{12}$ de sa valeur chaque année.
\item Les températures relevées à midi quatre jours de suite sont $21$, $23$, $22$ et $25$ degrés.
\item Un stock de $900$ articles diminue de $45$ articles par semaine.
\item Les effectifs d'un club, d'année en année, sont $40$, $44$, $49$ et $55$ adhérents.
\item Un capital de $\num{2000}$ euros est placé à $\pc{2,5}$ par an, intérêts composés.
\end{enumerate}

\end{activity}
\begin{solution}Une variation absolue constante indique le modèle arithmétique ; un pourcentage ou coefficient constant indique le modèle géométrique ; sinon aucun des deux n'est garanti. Classement attendu. Arithmétique : a) $r=\num{20000}$, $u_n=\num{600000}+\num{20000}n$ ; c) $r=2$, $u_n=15+2n$ ; f) $r=-45$, $u_n=900-45n$. Géométrique : b) $q=\num{1,03}$, $u_n=\num{5000}\times\num{1,03}^n$ ; d) $q=\num{0,88}$, $u_n=u_0\times\num{0,88}^n$ ; h) $q=\num{1,025}$, $u_n=\num{2000}\times\num{1,025}^n$. Ni l'une ni l'autre : e) les différences successives valent $2$, $-1$, $3$ et ne sont pas constantes, les quotients ne le sont pas non plus ; g) les différences valent $4$, $5$, $6$ et les quotients $\dfrac{44}{40}=\num{1,1}$, $\dfrac{49}{44}\approx\num{1,1136}$, $\dfrac{55}{49}\approx\num{1,1224}$ : aucun n'est constant. Faire verbaliser la donnée qui déclenche le choix de méthode.\end{solution}

\begin{summary}Avant de calculer, identifier : type de suite, premier rang, raison, terme ou somme demandé. La conclusion précise l'unité et le rang.\end{summary}

\begin{example}
Reprendre les données de l'activité et rédiger la solution dans l'ordre : définition des objets, propriété utilisée, calcul exact ou approché, puis phrase de réponse. La réponse modèle figure dans l'édition enseignant.
\end{example}
\begin{solution}Une variation absolue constante indique le modèle arithmétique ; un pourcentage ou coefficient constant indique le modèle géométrique ; sinon aucun des deux n'est garanti. Classement attendu. Arithmétique : a) $r=\num{20000}$, $u_n=\num{600000}+\num{20000}n$ ; c) $r=2$, $u_n=15+2n$ ; f) $r=-45$, $u_n=900-45n$. Géométrique : b) $q=\num{1,03}$, $u_n=\num{5000}\times\num{1,03}^n$ ; d) $q=\num{0,88}$, $u_n=u_0\times\num{0,88}^n$ ; h) $q=\num{1,025}$, $u_n=\num{2000}\times\num{1,025}^n$. Ni l'une ni l'autre : e) les différences successives valent $2$, $-1$, $3$ et ne sont pas constantes, les quotients ne le sont pas non plus ; g) les différences valent $4$, $5$, $6$ et les quotients $\dfrac{44}{40}=\num{1,1}$, $\dfrac{49}{44}\approx\num{1,1136}$, $\dfrac{55}{49}\approx\num{1,1224}$ : aucun n'est constant.\end{solution}

\begin{exercise}[essentiel]
\begin{enumerate}\item La suite $(u_n)$ vérifie $u_0=160$ et $u_{n+1}=u_n+3$. Écrire $u_n$ puis calculer $u_{12}$.\item Une quantité vaut $520$ et baisse de $\pc{3}$ par période. Écrire $v_n$ et calculer $v_5$.\item Pour chacune des suites $(u_n)$ et $(v_n)$ ci-dessus, préciser si le modèle est arithmétique ou géométrique et justifier par une différence ou un quotient.\end{enumerate}
\end{exercise}
\begin{solution}$u_n=160+3n$ et $u_{12}=196$. Pour la seconde, $q=\num{0,97}$, $v_n=520\times\num{0,97}^n$ et $v_5\approx\num{446,54}$. La suite $(u_n)$ est arithmétique (différence constante $3$) et $(v_n)$ est géométrique (quotient constant $\num{0,97}$).\end{solution}

\mxexuse{6}

\begin{homework}Reprendre la question 2 de l'exercice essentiel de la séance avec les données suivantes : remplacer la donnée initiale par $150$ et le taux ou la raison par $2$. Rédiger en quatre lignes : données, formule, calcul, conclusion.\end{homework}
\begin{solution}La correction dépend de l'exercice repris. Vérifier surtout la cohérence du modèle, le nombre d'itérations, l'arrondi et la phrase de conclusion.\end{solution}

\begin{exitticket}[3 minutes]Sur une demi-feuille : écrire la règle principale de la séance, traiter un exemple très court et indiquer son degré de confiance de 1 à 4. Question(s) essentielle(s) du jour : 1 et 2.\end{exitticket}

\ifmxteacher
\begin{method}
\textbf{Conduite en 1 h :} 7 min d'automatismes ; 10 min de mise en commun et trace ; 35 min d'exercices essentiels ; 5 min de correction ; 3 min de billet de sortie.\\
\textbf{Conduite en 2 h :} ajouter 20 min de recherche de l'activité, 20 min d'institutionnalisation, un second exercice autonome et 10 min de correction comparée.\\
\textbf{Différenciation :} fournir aux élèves fragiles une fiche à quatre cases (données, modèle, calcul, conclusion). Demander aux élèves rapides de produire un contre-exemple, une justification ou une seconde méthode.
\end{method}
\fi
